\section{Core Join Architecture and Execution Paths}
\label{sec:core-architecture}

This section presents the algorithmic core implemented in the current \system codebase. We focus on three designs that are jointly necessary for streaming vector similarity joins on multicore CPUs: (i) read/write-decoupled two-tier indexing, (ii) boundary-aware partition routing, and (iii) logical-partition remapping for skew. Together, they target high throughput, low latency, and robust recall under sliding-window updates.

\subsection{Execution Objective and Contract of \system}

\system operates on two unbounded streams with eager processing. For each incoming vector $x$ from side $s$, \system executes four stages in order: state update, candidate retrieval, verification, and emission. The contract is:
\begin{enumerate}
    \item \textbf{Low-overhead update}: ingest $x$ with minimal synchronization on the write path.
    \item \textbf{High-coverage retrieval}: collect candidates from complementary index views.
    \item \textbf{Exact verification}: enforce temporal and similarity predicates before output.
    \item \textbf{Bounded staleness}: refresh read-optimized structures asynchronously.
\end{enumerate}
This contract explicitly separates algorithmic approximation (candidate stage) from semantic correctness (verification stage).

Let an arriving record be $x$ from stream side $s \in \{L, R\}$ and let $\bar{s}$ denote the opposite side. \system emits $(x,y)$ only if
\[
|x.\mathit{ts} - y.\mathit{ts}| \leq W \quad \land \quad \mathrm{sim}(x,y) \geq \theta.
\]
Thus, recall loss can only come from candidate under-coverage, not from output semantics.

\subsection{Mechanism I: Read/Write-Decoupled Two-Tier Indexing}

The first key idea is decoupling update-intensive and query-intensive responsibilities:
\begin{itemize}
    \item \textbf{Local mutable indexes} absorb high-frequency inserts with partition-local ownership.
    \item \textbf{Global read-oriented indexes} provide broad recall-oriented probing.
\end{itemize}

For each query record $x$, \system probes both sources and merges candidates by identity:
\[
\mathcal{C}(x) = \mathrm{Dedup}(\mathcal{C}_{\text{local}}(x) \cup \mathcal{C}_{\text{global}}(x)).
\]
This avoids forcing one index to optimize conflicting goals (fast write vs. broad recall).

The local tier is optimized for update throughput and short critical sections. The global tier is optimized for query coverage and is refreshed in the background from window snapshots. Let $\Delta_r$ be rebuild interval and $\Delta_s$ be snapshot-to-query lag; then candidate freshness is bounded by approximately $\Delta_r + \Delta_s$, giving a controllable staleness knob rather than unpredictable drift.

To keep this bound meaningful under sparse or bursty arrivals, the current implementation computes rebuild validity against the latest event-time observed in snapshots (instead of raw wall-clock time). Concretely, if $t_{\max}$ is the maximum timestamp across the sampled left/right snapshots, global rebuild uses $t_{\max}-W$ as the effective lower bound and skips rebuild when both sides have no valid records. This avoids repeatedly rebuilding empty or semantically stale global views.

At high ingest rate, single-tier shared indexes suffer from either lock contention (if fully mutable) or stale views (if infrequently rebuilt). The two-tier design moves this trade-off from an implementation artifact to an explicit algorithm parameterization, where update pressure is absorbed locally and global recall is recovered periodically.

From a cost view, per-record hot-path work is
\[
T_{\text{hot}}(x) \approx T_{\text{route}}(x) + T_{\text{insert-local}}(x) + T_{\text{probe}}(x) + T_{\text{verify}}(x),
\]
while global rebuild cost is amortized asynchronously. This separation is central for keeping tail latency stable under bursty arrivals.

\subsection{Mechanism II: Boundary-Aware Partition Routing}

The second key idea is similarity-aware routing with controlled multicast. A strict single-partition policy risks missing near-border matches; unrestricted broadcast wastes compute and increases duplicates. In the current implementation, routing is realized in two stages:
\begin{itemize}
    \item \textbf{Connection-layer routing}: upstream-to-downstream queue selection through a partitioner (unicast or multicast).
    \item \textbf{VSJoin runtime routing}: records are mapped to logical partitions and then resolved to physical subtasks before local insert/query.
\end{itemize}
For VSJoin, \system uses boundary-aware fanout to navigate the recall/overhead trade-off:
\begin{itemize}
    \item Non-boundary vectors are routed to one primary partition.
    \item Boundary vectors are multicast to a small set of nearby partitions.
\end{itemize}

Let $\mathcal{P}(x)$ be the target partition set for record $x$. \system enforces:
\[
\mathbb{E}[|\mathcal{P}(x)|] \ll P,
\]
where $P$ is physical parallelism, while ensuring boundary vectors are replicated enough to prevent systematic misses. In practice, this is controlled by small fanout parameters (e.g., top-$k$ neighbor partitions or threshold-based overlap), and deduplication is applied before emission.

This mechanism addresses a core failure mode of naive partitioned vector joins: recall collapse as parallelism increases. Without boundary-aware fanout, partition granularity improves throughput but can sharply reduce cross-boundary match coverage.

In the implemented VSJoin path, boundary-aware routing is realized as a two-step mapping. First, an LSH partitioner outputs one or multiple physical partition candidates (multicast only for boundary vectors). Second, each selected physical partition is expanded to a logical partition with virtualization:
\[
\texttt{logical\_pid} = \texttt{physical\_pid} \times V + h(\texttt{uid}) \bmod V,
\]
where $V$ is virtual nodes per physical partition and $h(\cdot)$ is a stable hash. This design keeps routing deterministic for a record identity while enabling finer-grained balancing via logical remapping.

\subsection{Mechanism III: Logical-Partition Remapping for Skew}

The third key idea is decoupling partition semantics from worker placement. With runtime parallelism $P$ and virtualization factor $V$, \system exposes $P\times V$ logical partitions and maps them to physical workers through an assignment table. This indirection supports lightweight online rebalancing under skewed arrivals.

In the current code, logical partition id is computed as
\[
\texttt{logical\_pid} = \texttt{physical\_pid} \times V + \texttt{v\_idx},
\]
where $\texttt{v\_idx}$ is derived from a stable hash of record identity. Routing then applies
\[
\texttt{subtask} = \texttt{AssignmentTable}[\texttt{logical\_pid}].
\]

If worker load is $\ell_i$ for worker $i$, remapping seeks to reduce imbalance while preserving partition-local processing semantics. In practice, this improves robustness to hot partitions and non-uniform vector distributions, where static partition-to-thread binding often leads to persistent bottlenecks.

Concretely, \system maintains a lightweight load monitor that aggregates per-subtask samples (record volume, smoothed processing latency, and backlog signal when available). At each background tick, interval load is estimated from newly observed records (plus a small backlog weight), and rebalance is triggered when
\[
\frac{\max_i \hat{\ell}_i}{\frac{1}{P}\sum_i \hat{\ell}_i} > \tau,
\]
where $\tau$ is a configurable imbalance ratio. When triggered, up to $M$ logical partitions are moved from the busiest subtask to the idlest one in one atomic AssignmentTable update ($M$ is also configurable).

More specifically, if $R_i^{(k)}$ is cumulative reported records for worker $i$ at tick $k$, and $B_i^{(k)}$ is the EWMA backlog estimate, the interval load signal is:
\[
\hat{\ell}_i^{(k)} = \big(R_i^{(k)}-R_i^{(k-1)}\big) + \beta B_i^{(k)},
\]
with a small constant $\beta$ (backlog weight). The migration count in one round is bounded by
\[
m = \min\big(M,\; \text{overload\_partitions},\; \text{movable\_candidates}\big),
\]
which guarantees predictable remapping overhead.

The AssignmentTable uses an RCU-style double-buffer swap: reads are lock-free, while updates are batched and atomically published.

The table update path is fail-safe in the current code: invalid logical/physical sizes disable remapping gracefully, and each update applies to a copied snapshot before atomic pointer publication. This keeps concurrent readers wait-free and avoids transient partial mappings.

The present rebalancer is intentionally lightweight and heuristic: it uses sampled per-subtask load signals and caps migration count per round. Therefore, it improves robustness under skew but does not claim globally optimal balancing each round.

In the current implementation, remapping updates routing metadata only (i.e., future routing decisions) and avoids synchronous full-state migration on the foreground path. This design keeps rebalance overhead predictable under high ingest pressure.

This mechanism is essential for real streams where data distribution drifts over time. A fixed mapping can be optimal at startup but degrade later; remapping keeps throughput more stable over time without changing query semantics.

\subsection{End-to-End Pipeline in \system}

Combining the three mechanisms yields the following pipeline for each record: (i) route with boundary-aware fanout to one or multiple logical partitions, (ii) map logical partitions to physical subtasks via AssignmentTable, (iii) insert into local mutable index/state at target subtasks, (iv) retrieve from local+global candidate sources, (v) deduplicate and verify, and (vi) emit exact-threshold join pairs. Concurrently, a background task rebuilds global indexes from fresh window snapshots and triggers periodic rebalance checks.

The current implementation executes these background actions in a single periodic control loop: each tick first evaluates skew and optionally remaps logical partitions, then performs global-index refresh using event-time-valid snapshots. This ordering prioritizes routing adaptation before rebuilding read-optimized structures, so subsequent rebuilds better reflect updated placement.

This pipeline is the core algorithmic contribution of \system: it turns the classical trade-off between shared-index contention and partition-induced recall loss into a controllable design space using decoupled indexing, bounded multicast, and logical remapping.

From a systems perspective, the current implementation provides three concrete guarantees: (1) semantic correctness is enforced at verification time; (2) routing-map updates are atomic for concurrent readers; and (3) rebalance overhead is bounded by per-round migration limits. At the same time, we do not assume perfect balance under all skews; observed balance quality depends on load signal quality, rebalance interval, and workload drift speed.

\subsection{Role of Baselines}

BruteForce, IVF, HNSW, HDR-Tree, ClusteredJoin, and S3J-style paths are included as comparison targets under the same execution semantics. They are not the narrative center of this paper; instead, they quantify how each \system key design contributes to quality and performance in the same workload and hardware setting.

The next section presents implementation and evaluation details used to measure these effects.
